\section{Consistency and the value of auxiliary records}

The risk comparison in \cref{obj:thm:uniform-annotation-frontier} answers three resource questions: which budgets support consistency, how much marginal information attains rare-label oracle order, and when auxiliary records yield a vanishing risk ratio relative to complete records alone. These questions allow the sample sizes, alphabet size, and overlap floor to vary jointly. We first specify the domain of experiment sequences.

\begin{definitionv}{synth_8}[Allowed experiment sequences]
An experiment sequence \(\mathbf v\) is an arbitrary sequence of public experiment indices, indexed by the nonnegative integers:
\[
\mathbf v=\bigl((n_k,m_k,d_k,\epsilon_k)\bigr)_{k=0}^{\infty}.
\]
For every \(k\), the complete-record size \(n_k\), auxiliary size \(m_k\), and alphabet size \(d_k\) are integers satisfying
\[
n_k\ge 1,\qquad m_k\ge 0,\qquad d_k\ge 2,
\]
and the public overlap floor \(\epsilon_k\) is a real number satisfying
\[
0<\epsilon_k\le \frac14.
\]
\end{definitionv}

An experiment sequence \leanref{sym:\mathbf v}{\ensuremath{\mathbf v}} permits arbitrary joint variation within these public ranges, including diminishing overlap and an expanding alphabet. The three resource criteria concern vanishing minimax risk, bounded risk relative to the rare-label benchmark, and vanishing risk relative to the supervised experiment. The benchmark retains its cap:
\[
b(n,\epsilon)=\min\{1,(n\epsilon)^{-1}\}.
\]
The following definition makes these comparisons precise and introduces the sequence-specific information scales.

\begin{definitionv}{def:phases}[Resource phases \(\mathcal C,\mathcal O,\mathcal I\)]
Define
\[
\begin{aligned}
\mathcal C
&=\left\{\mathbf v:
R(n_k,m_k,d_k,\epsilon_k)\longrightarrow 0\right\},\\
\mathcal O
&=\left\{\mathbf v:
n_k\epsilon_k\longrightarrow\infty,\quad
\limsup_{k\to\infty}
\frac{R(n_k,m_k,d_k,\epsilon_k)}{b(n_k,\epsilon_k)}
<\infty\right\},\\
\mathcal I
&=\left\{\mathbf v:
\frac{R(n_k,m_k,d_k,\epsilon_k)}
{R(n_k,0,d_k,\epsilon_k)}
\longrightarrow 0\right\}.
\end{aligned}
\]
Each set ranges over all allowed sequences of public experiment indices
\[
\mathbf v=((n_k,m_k,d_k,\epsilon_k))_{k\ge1}.
\]
For each such sequence, the rare-arm information scale and total marginal-information sample size are
\[
S_k=n_k\epsilon_k,\qquad N_k=n_k+m_k.
\]
Write \(\ell_k\) for the regularized logarithm \(\ell\) evaluated at the \(k\)-th experiment.
\end{definitionv}

Explicitly, the regularized logarithm is \(\ell_k=\log(\exp(1)+S_k)\). Sequences are indexed by \(k\ge0\) as in \cref{obj:synth_8}; the \(k\ge1\) display in \cref{obj:def:phases} denotes the same asymptotic tail and discards only the initial coordinate. The consistency phase \leanref{sym:\mathcal C}{\ensuremath{\mathcal C}} concerns absolute precision. The oracle phase \leanref{sym:\mathcal O}{\ensuremath{\mathcal O}} combines a diverging rare-arm information scale with attainment of the rare-label benchmark up to a constant factor. The improvement phase \leanref{sym:\mathcal I}{\ensuremath{\mathcal I}} measures a strict relative gain from auxiliary records. The next result characterizes all three criteria and supplies finite-experiment conditions for saturation and rare-label order.

\begin{theoremv}{thm:annotation-resource-phases}[Exact annotation resource phases]
Let \(\mathbf v=((n_k,m_k,d_k,\epsilon_k))_{k\ge0}\) be any sequence of allowed public experiment indices, and put
\[
S_k=n_k\epsilon_k,\qquad
N_k=n_k+m_k,\qquad
\ell_k=\log(\exp(1)+S_k).
\]
For the phases defined in \cref{obj:def:phases}, the following equivalences hold as \(k\to\infty\):
\[
\begin{aligned}
\mathbf v\in\mathcal C
&\Longleftrightarrow
S_k\longrightarrow\infty
\quad\text{and}\quad
d_k=o(N_k\epsilon_k\ell_k),\\
\mathbf v\in\mathcal O
&\Longleftrightarrow
S_k\longrightarrow\infty
\quad\text{and}\quad
d_k=O\!\left(\frac{N_k\epsilon_k\ell_k}{\sqrt{S_k}}\right),\\
\mathbf v\in\mathcal I
&\Longleftrightarrow
S_k\longrightarrow\infty,\quad
\frac{d_k}{\sqrt{S_k}\ell_k}\longrightarrow\infty,\quad
\frac{N_k/n_k}{\max\{1,d_k/(S_k\ell_k)\}}\longrightarrow\infty.
\end{aligned}
\]
If \(m_k=0\) for every \(k\), then
\[
\begin{aligned}
\mathbf v\in\mathcal C
&\Longleftrightarrow
S_k\longrightarrow\infty
\quad\text{and}\quad d_k=o(S_k\ell_k),\\
\mathbf v\in\mathcal O
&\Longleftrightarrow
S_k\longrightarrow\infty
\quad\text{and}\quad d_k=O(\sqrt{S_k}\ell_k),
\end{aligned}
\qquad
\mathbf v\notin\mathcal I.
\]
Whenever \(S_k\to\infty\), oracle membership also has the equivalent budget characterization
\[
\mathbf v\in\mathcal O
\quad\Longleftrightarrow\quad
\frac{d_k\sqrt{S_k}}{\epsilon_k\ell_k}=O(N_k).
\]

Write \(R\) for the minimax risk in \cref{obj:def:risk}. For every \(\mathbf v\in\mathcal O\), there exist constants \(c,C>0\) such that, for all sufficiently large \(k\),
\[
\frac{c}{S_k}
\le R(n_k,m_k,d_k,\epsilon_k)
\le \frac{C}{S_k}.
\]
Moreover, there exists \(c>0\) such that, for all sufficiently large \(k\) and every \(m'\in\mathbb N\),
\[
\frac{c}{S_k}\le R(n_k,m',d_k,\epsilon_k).
\]

For a single experiment, let \(n,m,d\in\mathbb N\) and \(\epsilon\in\mathbb R\) satisfy
\begin{itemize}
\item \textbf{(Complete-record budget.)} \(n\ge1\).
\item \textbf{(Alphabet size.)} \(d\ge2\).
\item \textbf{(Overlap range.)} \(0<\epsilon\le1/4\).
\end{itemize}
Put
\[
N=n+m,\qquad S=n\epsilon,\qquad
\ell=\log(\exp(1)+S),\qquad
b(n,\epsilon)=\min\{1,S^{-1}\}.
\]
There exists a universal constant \(c>0\) such that, for every such experiment,
\[
N\epsilon\ell\le d
\quad\Longrightarrow\quad
c\le R(n,m,d,\epsilon)\le1.
\]
There also exist universal constants \(c,C>0\) such that, for every such experiment,
\[
\frac{d\sqrt{n\epsilon}}{\epsilon\ell}\le N
\quad\Longrightarrow\quad
c\,b(n,\epsilon)\le R(n,m,d,\epsilon)\le C\,b(n,\epsilon).
\]
\end{theoremv}

The theorem separates the two resources directly. Complete records determine \(S_k\), so no auxiliary budget can overcome bounded rare-arm outcome information. Both channels determine \(N_k\), so auxiliary records can reduce the many-cell term once \(S_k\to\infty\). For a single experiment, the two useful endpoints are
\[
N\epsilon\ell\le d
\quad\Longrightarrow\quad R\asymp1,
\qquad
N\ge \frac{d\sqrt S}{\epsilon\ell}
\quad\Longrightarrow\quad
R\asymp \min\{1,S^{-1}\},
\]
with universal comparison constants. The first is marginal-information saturation; the second is the rare-label plateau.

\subsection{Three illustrative sequences}

Let \(t=k+4\), so every displayed public index is legal and \(\ell_k\asymp\log t\).
\begin{enumerate}
\item \textbf{Supervised rare-label order.} Take
\[
n_k=t^4,\qquad \epsilon_k=t^{-1},\qquad m_k=0,\qquad d_k=t.
\]
Then \(S_k=t^3\) and \(d_k/(\sqrt{S_k}\ell_k)\to0\). The supervised experiment is already in \(\mathcal O\), and \(R(n_k,0,d_k,\epsilon_k)\asymp t^{-3}\).

\item \textbf{Improvement between two consistent risks.} Take
\[
n_k=t^4,\qquad \epsilon_k=t^{-1},\qquad d_k=t^2,\qquad m_k=t^6.
\]
Without auxiliary records the many-cell term gives
\[
R(n_k,0,d_k,\epsilon_k)\asymp \frac{1}{t^2(\log t)^2},
\]
while the added marginal budget reduces the risk to \(R(n_k,m_k,d_k,\epsilon_k)\asymp t^{-3}\). Both risks vanish, but their ratio is of order \((\log t)^2/t\to0\), so the sequence lies in \(\mathcal I\).

\item \textbf{Expansion from saturation to consistency.} Take
\[
n_k=t^4,\qquad \epsilon_k=t^{-1},\qquad d_k=t^5,\qquad m_k=t^8.
\]
The supervised scale satisfies \(n_k\epsilon_k\ell_k\asymp t^3\log t\ll d_k\), so its risk stays of constant order. With auxiliary records,
\((n_k+m_k)\epsilon_k\ell_k\asymp t^7\log t\), making
\(d_k=o((n_k+m_k)\epsilon_k\ell_k)\); the two-channel risk is then of order \(t^{-3}\) and is consistent.
\end{enumerate}

These examples are illustrations of the order characterizations, rather than finite-budget calibrations. They also show why strict improvement can compare two vanishing risks or move an experiment out of saturation, while the rare-label plateau remains controlled by complete-record outcomes.
